\documentclass[11pt]{article}
\usepackage[T1]{fontenc}
\usepackage{lmodern,microtype}
\usepackage[margin=0.95in]{geometry}
\usepackage{amsmath,amssymb,amsthm,booktabs,array}
\usepackage{fancyhdr}
\usepackage[colorlinks=true,linkcolor=blue,urlcolor=blue,citecolor=blue]{hyperref}
\usepackage{enumitem}
\setlength{\parindent}{0pt}
\setlength{\parskip}{0.55em}
\setlength{\headheight}{14pt}
\emergencystretch=2em
\pagestyle{fancy}
\fancyhf{}
\fancyhead[L]{\small Knights and rooks: finite nonexistence}
\fancyhead[R]{\small Research record}
\fancyfoot[C]{\thepage}
\newtheorem{theorem}{Theorem}
\newtheorem{lemma}{Lemma}
\newcommand{\Z}{\mathbb Z}
\newcommand{\Kset}{\mathcal K}
\newcommand{\Rset}{\mathcal R}
\hypersetup{pdftitle={Nonexistence in the N=4, R=2 Knights-and-Rooks Problem},pdfauthor={Research record prepared in this conversation}}
\title{\vspace{-1em}Nonexistence in the \(N=4,\ R=2\)\\Knights-and-Rooks Problem}
\author{Self-contained proof and independently checkable certificate}
\date{13 September 2026}
\begin{document}
\maketitle
\thispagestyle{fancy}
\begin{abstract}
There is no finite placement in which every knight attacks exactly four
rooks and no knights, while every rook sees exactly two knights and no
rooks. A rightmost-knight argument proves a stronger result: if every
rook sees exactly two knights and no rooks, then some knight attacks at
most three rooks, even when knight--knight attacks are permitted. A
separately encoded necessary local relaxation is also proved
unsatisfiable by a small, independently checked certificate. The universal
conclusion does not rely on exhausting boards of bounded size.
\end{abstract}

\section{Exact model and result}
Let \(\Kset,\Rset\subset\Z^2\) be finite, nonempty, disjoint sets. Their elements
are knight and rook squares, respectively. A knight at \((x,y)\) attacks
occupied squares whose absolute coordinate differences from \((x,y)\)
are \(1\) and \(2\). It jumps over intervening pieces.

For a rook, inspect each of the four open orthogonal rays and retain only
its first occupied square, when one exists. The rook condition is that
exactly two such first squares contain knights and none contain rooks.
Consequently, exactly two rays are occupied and the other two rays are
\emph{entirely empty}. Attacks are directed; no balance between the two
kinds of outgoing attack counts is assumed.

\begin{theorem}\label{thm:main}
If every rook sees exactly two knights and no rooks, then some knight
attacks at most three rooks. This holds without a prohibition on
knight--knight attacks.
\end{theorem}
The requested placement would give every knight four rook targets, so
Theorem~\ref{thm:main} rules it out.

\section{Universal geometric proof}
Assume for contradiction that every knight attacks at least four rooks.

\paragraph{A rightmost knight bounds all pieces.}
Every rook has \(x\)-coordinate at most the maximum knight \(x\)-coordinate.
Indeed, a rook strictly to the right of all knights could see a knight
only to its west, giving at most one knight target. Translate a rightmost
knight to
\[
A=(0,0).
\]
Every square with \(x>0\) is therefore empty. The finiteness assumption is
used here to choose a rightmost knight, not to impose a board-size bound.

\paragraph{Four forced rooks and their blockers.}
Only four knight-move destinations of \(A\) have \(x\leq0\). All four must
be rooks:
\[
U=(-2,1),\qquad V=(-1,2),\qquad
W=(-1,-2),\qquad Z=(-2,-1).
\]
The rooks \(U,Z\) must have a knight between them; their sole intervening
square is
\[
C=(-2,0),
\]
which is consequently a knight. Similarly, \(V,W\) require a knight at
at least one of
\[
B_+=(-1,1),\qquad B_0=(-1,0),\qquad B_-=(-1,-1).
\]

\paragraph{The middle possibility is impossible.}
If \(B_0\) were a knight, its two knight-move destinations with
\(x=1\) would be empty. Each of its other six destinations is immediately
horizontally adjacent to a forced rook:
\[
\begin{array}{c|c}
\text{destination from }B_0&\text{adjacent forced rook}\\\hline
(-3,1),\ (-3,-1)&U,\ Z\\
(-2,2),\ (0,2)&V\\
(-2,-2),\ (0,-2)&W
\end{array}
\]
None of these six destinations can be a rook: two horizontally adjacent
rooks would see each other. Thus a knight at \(B_0\) would attack no
rooks, a contradiction. At least one of \(B_+,B_-\) is a knight. Reflect
in \(y=0\), if needed, so that \(B_+\) is a knight. This reflection resolves
a proved two-case alternative; it is not an assumed symmetry of the placement.

\paragraph{The upper blocker forces a right-edge rook.}
The rook \(U\) already sees \(C\) immediately south and \(B_+\) immediately
east. Its north and west rays must be entirely empty. In particular,
\((-2,2)\) and \((-2,3)\) are empty. Also \((-3,2)\) cannot be a rook,
because it would see \(V=(-1,2)\) across empty \((-2,2)\).

Of the eight knight-move destinations of \(B_+\), only
\[
(-3,0),\quad(-2,-1),\quad(0,-1),\quad(0,3)
\]
can therefore contain rooks. Two other destinations have \(x>0\), and
the remaining two were excluded above. All four listed squares must be
rooks. In particular,
\[
X=(0,-1)
\]
is a rook. The rooks \(Z,X\) have only one intervening square, so
\(B_-=(-1,-1)\) is now forced to be a knight.

\paragraph{The lower blocker has at most three targets.}
The rook \(Z\) sees \(C\) immediately north and \(B_-\) immediately east,
so its south ray is empty. In particular, \((-2,-2)\) and \((-2,-3)\)
are empty. Therefore \((-3,-2)\) cannot be a rook, because it would see
\(W=(-1,-2)\) across empty \((-2,-2)\).

The rook \(X\) sees \(A\) immediately north and \(B_-\) immediately west.
Its south ray is empty, so \((0,-3)\) is empty as well. The full list of
knight-move destinations of \(B_-\) is
\[
\begin{split}
&(-3,-2),\ (-3,0),\ (-2,-3),\ (-2,1),\\
&(0,-3),\ (0,1),\ (1,-2),\ (1,0).
\end{split}
\]
The last two are in the empty half-plane \(x>0\), and \((-3,-2)\),
\((-2,-3)\), \((0,-3)\) were just excluded. At most three rook targets
remain:
\[
(-3,0),\qquad(-2,1),\qquad(0,1).
\]
This contradicts the lower bound of four for the knight \(B_-\), proving
the theorem.\hfill\(\square\)

No step uses the absence of knight--knight attacks, a fixed piece inventory,
board colors, legal-game reachability, or reciprocity between attacks.

\section{A separately checkable finite obstruction}
The proof above is sufficient on its own. The following computation was
used to try to falsify its local reasoning and to provide a checkable
certificate.

\subsection{Why this is not bounded-board exhaustion}
After the proved rightmost-knight reduction, retain only the window
\[
D=\{-3,-2,-1,0\}\times\{-3,-2,-1,0,1,2,3\}
\]
and the four distinguished positions
\[
T=\{(0,0),(-1,0),(-1,1),(-1,-1)\}.
\]
Every possible rook target of a knight at a position in \(T\) lies either
in \(D\) or in the already-proved empty half-plane \(x>0\). Consequently,
their knight-degree requirements can be checked locally without assuming
anything about other knights. Pieces outside \(D\) are otherwise permitted.

A local rook has at most two locally occupied rays: outside pieces cannot
remove a ray that already contains a local occupied square. Every aligned
pair of rooks in \(D\) must have a knight between them; all intervening
squares lie in \(D\). Thus any hypothetical finite counterexample induces
an assignment satisfying the following \emph{necessary relaxation}.

\subsection{Complete encoding}
For each \(p\in D\), let \(K_p,R_p,O_p\) be Boolean variables for knight,
rook, and occupancy. For every nonempty intersection of an open orthogonal
ray from \(p\) with \(D\), let \(V_{p,d}\) mean that this intersection is
occupied. The constraints are
\begin{align*}
&\neg(K_p\wedge R_p),\qquad O_p\leftrightarrow(K_p\vee R_p),\\
&V_{p,d}\leftrightarrow\bigvee_{q\in D\cap\operatorname{ray}(p,d)} O_q,\\
&R_p\ \Longrightarrow\ \sum_d V_{p,d}\leq2,\\
&R_p\wedge R_q\ \Longrightarrow\ \bigvee_{s\text{ strictly between }p,q}K_s
   \quad\text{for aligned }p\ne q,\\
&K_{(0,0)},\\
&K_p\ \Longrightarrow\ \sum_{q\in D:\ q-p\text{ is a knight move}} R_q\geq4
   \quad(p\in T).
\end{align*}
A ray with empty intersection contributes the constant false. There are
no degree constraints on other knights, no knight--knight restriction,
and no lower bound of two on the \emph{local} rook ray count. These omissions
relax the target, rather than exclude possible counterexamples.

The generator expands the equivalences into implications, uses one guarded
clause for each triple of ray flags for the at-most-two constraints, and
encodes at least four among \(t\) target variables by requiring every
\((t-3)\)-element subset to contain a true variable. Here \(t=4\) at the
origin and \(t=6\) at the other distinguished positions.

\begin{center}
\begin{tabular}{lr}
\toprule
Clause group&Count\\\midrule
Disjoint types and occupancy equivalences&112\\
Occupied-ray equivalences&342\\
Rook occupied-ray upper bounds&54\\
Intervening-knight conditions&126\\
Conditional knight lower bounds&64\\
Origin is a knight&1\\\midrule
Total&699\\\bottomrule
\end{tabular}
\end{center}
There are 174 Boolean variables: 84 type/occupancy variables and 90 ray
variables. The full DIMACS formula, variable map, and clause explanations
are included in the bundle.

\subsection{Certificate and verification}
The separate DPLL prover produced an UNSAT tree with three nodes: a single
split on \(K_{(-1,1)}\) and two contradictory leaves. Across the tree there
are 151 unit inferences. Every unit inference cites an original clause;
every leaf cites an original clause falsified by its branch assignment.

The independent checker \texttt{src/verify\_unsat.py} verifies the formula's
SHA-256 digest, each unit inference, both branches of the split, and both
conflicts. It imports neither the formula generator nor the prover, attack
checkers, or Z3. Its inspected output was
\begin{verbatim}
VERIFIED UNSAT {"conflicts": 2, "nodes": 3,
                "splits": 1, "unit_steps": 151}
\end{verbatim}
Independence here means separate algorithms and a proof checker that does
not trust the search code. It is not a claim of external peer review or
proof-assistant formalization. The human-readable reduction explains why
the local UNSAT certificate implies universal finite nonexistence.

\section{Calibration, falsification tests, and reproducibility}
Two independent attack reconstruction programs were executed. The first
uses pairwise knight-coordinate differences and selects the nearest
occupied point in each rook ray. The second generates the eight knight
jumps explicitly and obtains rook visibility from neighboring entries in
sorted occupied rows and columns.

The source's \(N=1,R=1\) and \(N=4,R=1\) illustrations were fetched,
visually transcribed, and accepted by both programs~\cite{n1r1,n4r1}.
Their original Git blob hashes were reproduced from the decoded image
bytes. The latter illustration has two knights and eight rooks; it is
not a solution to the \(N=4,R=2\) target. All its per-piece outgoing attacks
are retained in the accompanying JSON report.

The search also reproduced an \(N=2,R=3\) calibration with eight knights
and six rooks. Direct reconstruction yields 16 knight-to-rook attacks and
18 rook-to-knight attacks. This unequal count explicitly guards against
incorrectly imposing reciprocity.

Executed validation included 2,000 seeded random placements, eight
translated/reflected/rotated versions of a source example, invalid-piece
and blocker tests, five deliberately corrupted certificates (all rejected),
and three constraint-removal controls (all verified satisfiable). The
full suite ran 20 test methods successfully and regenerated a certificate
identical to the bundled one.

\paragraph{Earlier searches and their limits.}
Exploration used Python 3.13.5 and a locally installed Z3 4.13.3.0 shared
library. Unrestricted \(8\times8\) and \(16\times16\) target instances
reported UNSAT; symmetry-restricted \(12\times12\) and \(16\times16\)
instances also reported UNSAT. A \(32\times32\) attempt was interrupted
by the execution time limit and has no reported outcome. None of these
bounded searches proves the theorem. Their formulas, seeds, outputs, and
limitations are preserved separately.

\paragraph{Run the final verification.}
The final verification needs only Python 3.9 or newer and its standard
library. No network, pip package, or external solver is needed:
\begin{verbatim}
python3 src/run_all.py
\end{verbatim}
For the minimal certificate check alone:
\begin{verbatim}
python3 src/verify_unsat.py \
  certificates/local_relaxation.cnf \
  certificates/local_unsat_tree.json
\end{verbatim}
The bundle also contains \texttt{PROOF.md}, \texttt{FINITE\_REDUCTION.md},
source coordinates and attack reports, executable source code, inspected
logs, a checksum manifest, and \texttt{progress.json} with the exact
verification/resume command.

\section{Source conventions and historical status}
The target is column \(N=4\), row \(R=2\), in the Knights and Rooks table
of Friedman's February 2007 page~\cite{friedman}. Its accompanying text
specifies outgoing attacks on the other type only. The column and row
interpretations agree with the independently checked illustrations.
The target cell still displayed a question mark when rechecked after the
proof was found. A later puzzle post concerns the different \(N=4,R=1\)
case, not this target~\cite{puzzling}.

Initial and final web searches did not locate an earlier exact-target
solution or nonexistence proof. This is not a comprehensive proof of
historical priority, and an archival question mark alone does not establish
current open status. No priority claim, external review, publication, or
personal contact is asserted. The mathematical answer follows from the
explicit proof and certificate above, independently of that historical
question.

\begin{thebibliography}{9}
\bibitem{friedman} Erich Friedman. \emph{Problem of the Month (February 2007)},
Knights and Rooks table and problem statement.
\url{https://erich-friedman.github.io/mathmagic/0207.html}.
Accessed and rechecked in this research session, 13 September 2026.
\bibitem{n1r1} Erich Friedman. \(N=1,R=1\) source illustration.
\url{https://erich-friedman.github.io/mathmagic/0207/N1R1.gif}.
\bibitem{n4r1} Erich Friedman. \(N=4,R=1\) source illustration.
\url{https://erich-friedman.github.io/mathmagic/0207/N4R1.gif}.
\bibitem{puzzling} Will.Octagon.Gibson. \emph{Where should the chess pieces be
placed in 4x4, 6x6 and 9x9 chessboards?} Puzzling Stack Exchange,
26 August 2024; edited 29 May 2025. Puzzle 2 is \(N=4,R=1\).
\url{https://puzzling.stackexchange.com/questions/128060/}.
\end{thebibliography}
\end{document}
